\documentclass[CJK]{beamer}
\usepackage{CJKutf8}
\usepackage{beamerthemesplit}
\usetheme{Malmoe}
\useoutertheme[footline=authortitle]{miniframes}
\usepackage{amsmath}
\usepackage{amssymb}
\usepackage{graphicx}
\usepackage{eufrak}
\usepackage{color}
\usepackage{slashed}
\usepackage{simplewick}
\usepackage{tikz}
\usepackage{tcolorbox}
\graphicspath{{../figures/}}
%%figures
\def\lfig#1#2{\includegraphics[width=#1 in]{#2}}
\def\addfig#1#2{\begin{center}\includegraphics[width=#1 in]{#2}\end{center}}
\def\wulian{\includegraphics[width=0.18in]{emoji_wulian.jpg}}
\def\bigwulian{\includegraphics[width=0.35in]{emoji_wulian.jpg}}
\def\bye{\includegraphics[width=0.18in]{emoji_bye.jpg}}
\def\bigbye{\includegraphics[width=0.35in]{emoji_bye.jpg}}
\def\huaixiao{\includegraphics[width=0.18in]{emoji_huaixiao.jpg}}
\def\bighuaixiao{\includegraphics[width=0.35in]{emoji_huaixiao.jpg}}
%% colors
\def\blacktext#1{{\color{black}#1}}
\def\bluetext#1{{\color{blue}#1}}
\def\redtext#1{{\color{red}#1}}
\def\darkbluetext#1{{\color[rgb]{0,0.2,0.6}#1}}
\def\skybluetext#1{{\color[rgb]{0.2,0.7,1.}#1}}
\def\cyantext#1{{\color[rgb]{0.,0.5,0.5}#1}}
\def\greentext#1{{\color[rgb]{0,0.7,0.1}#1}}
\def\darkgray{\color[rgb]{0.2,0.2,0.2}}
\def\lightgray{\color[rgb]{0.6,0.6,0.6}}
\def\gray{\color[rgb]{0.4,0.4,0.4}}
\def\blue{\color{blue}}
\def\red{\color{red}}
\def\green{\color{green}}
\def\darkgreen{\color[rgb]{0,0.4,0.1}}
\def\darkblue{\color[rgb]{0,0.2,0.6}}
\def\skyblue{\color[rgb]{0.2,0.7,1.}}
%%control
\def\be{\begin{equation}}
\def\ee{\nonumber\end{equation}}
\def\bea{\begin{eqnarray}}
\def\eea{\nonumber\end{eqnarray}}
\def\bch{\begin{CJK}{UTF8}{gbsn}}
\def\ech{\end{CJK}}
\def\bitem{\begin{itemize}}
\def\eitem{\end{itemize}}
\def\bcenter{\begin{center}}
\def\ecenter{\end{center}}
\def\bex{\begin{minipage}{0.2\textwidth}\includegraphics[width=0.6in]{jugelizi.png}\end{minipage}\begin{minipage}{0.76\textwidth}}
\def\eex{\end{minipage}}
\def\chtitle#1{\frametitle{\bch#1\ech}}
\def\bmat#1{\left(\begin{array}{#1}}
\def\emat{\end{array}\right)}
\def\bcase#1{\left\{\begin{array}{#1}}
\def\ecase{\end{array}\right.}
\def\bmini#1{\begin{minipage}{#1\textwidth}}
\def\emini{\end{minipage}}
\def\tbox#1{\begin{tcolorbox}#1\end{tcolorbox}}
\def\pfrac#1#2#3{\left(\frac{\partial #1}{\partial #2}\right)_{#3}}
%%symbols
\def\sone{$\star$}
\def\stwo{$\star\star$}
\def\sthree{$\star\star\star$}
\def\sfour{$\star\star\star\star$}
\def\sfive{$\star\star\star\star\star$}
\def\rint{{\int_\leftrightarrow}}
\def\roint{{\oint_\leftrightarrow}}
\def\stdHf{{\textit{\r H}_f}}
\def\deltaH{{\Delta \textit{\r H}}}
\def\ii{{\dot{\imath}}}
\def\skipline{{\vskip0.1in}}
\def\skiplines{{\vskip0.2in}}
\def\lagr{{\mathcal{L}}}
\def\hamil{{\mathcal{H}}}
\def\vecv{{\mathbf{v}}}
\def\vecx{{\mathbf{x}}}
\def\vecy{{\mathbf{y}}}
\def\veck{{\mathbf{k}}}
\def\vecp{{\mathbf{p}}}
\def\vecn{{\mathbf{n}}}
\def\vecA{{\mathbf{A}}}
\def\vecP{{\mathbf{P}}}
\def\vecsigma{{\mathbf{\sigma}}}
\def\hatJn{{\hat{J_\vecn}}}
\def\hatJx{{\hat{J_x}}}
\def\hatJy{{\hat{J_y}}}
\def\hatJz{{\hat{J_z}}}
\def\hatj#1{\hat{J_{#1}}}
\def\hatphi{{\hat{\phi}}}
\def\hatq{{\hat{q}}}
\def\hatpi{{\hat{\pi}}}
\def\vel{\upsilon}
\def\Dint{{\mathcal{D}}}
\def\adag{{\hat{a}^\dagger}}
\def\bdag{{\hat{b}^\dagger}}
\def\cdag{{\hat{c}^\dagger}}
\def\ddag{{\hat{d}^\dagger}}
\def\hata{{\hat{a}}}
\def\hatb{{\hat{b}}}
\def\hatc{{\hat{c}}}
\def\hatd{{\hat{d}}}
\def\hatN{{\hat{N}}}
\def\hatH{{\hat{H}}}
\def\hatp{{\hat{p}}}
\def\Fup{{F^{\mu\nu}}}
\def\Fdown{{F_{\mu\nu}}}
\def\newl{\nonumber \\}
\def\vece{\mathrm{e}}
\def\calM{{\mathcal{M}}}
\def\calT{{\mathcal{T}}}
\def\calR{{\mathcal{R}}}
\def\barpsi{\bar{\psi}}
\def\baru{\bar{u}}
\def\barv{\bar{\upsilon}}
\def\qeq{\stackrel{?}{=}}
\def\torder#1{\mathcal{T}\left(#1\right)}
\def\rorder#1{\mathcal{R}\left(#1\right)}
\def\contr#1#2{\contraction{}{#1}{}{#2}#1#2}
\def\trof#1{\mathrm{Tr}\left(#1\right)}
\def\trace{\mathrm{Tr}}
\def\comm#1{\ \ \ \left(\mathrm{used}\ #1\right)}
\def\tcomm#1{\ \ \ (\text{#1})}
\def\slp{\slashed{p}}
\def\slk{\slashed{k}}
\def\calp{{\mathfrak{p}}}
\def\veccalp{\mathbf{\mathfrak{p}}}
\def\Tthree{T_{\tiny \textcircled{3}}}
\def\pthree{p_{\tiny \textcircled{3}}}
\def\dbar{{\,\mathchar'26\mkern-12mu d}}
\def\erf{\mathrm{erf}}
\def\const{\mathrm{constant}}
\def\pheat{\pfrac p{\ln T}V}
\def\vheat{\pfrac V{\ln T}p}
%%units
\def\fdeg{{^\circ \mathrm{F}}}
\def\cdeg{^\circ \mathrm{C}}
\def\atm{\,\mathrm{atm}}
\def\angstrom{\,\text{\AA}}
\def\SIL{\,\mathrm{L}}
\def\SIkm{\,\mathrm{km}}
\def\SIyr{\,\mathrm{yr}}
\def\SIGyr{\,\mathrm{Gyr}}
\def\SIV{\,\mathrm{V}}
\def\SImV{\,\mathrm{mV}}
\def\SIeV{\,\mathrm{eV}}
\def\SIkeV{\,\mathrm{keV}}
\def\SIMeV{\,\mathrm{MeV}}
\def\SIGeV{\,\mathrm{GeV}}
\def\SIcal{\,\mathrm{cal}}
\def\SIkcal{\,\mathrm{kcal}}
\def\SImol{\,\mathrm{mol}}
\def\SIN{\,\mathrm{N}}
\def\SIHz{\,\mathrm{Hz}}
\def\SIm{\,\mathrm{m}}
\def\SIcm{\,\mathrm{cm}}
\def\SIfm{\,\mathrm{fm}}
\def\SImm{\,\mathrm{mm}}
\def\SInm{\,\mathrm{nm}}
\def\SImum{\,\mathrm{\mu m}}
\def\SIJ{\,\mathrm{J}}
\def\SIW{\,\mathrm{W}}
\def\SIkJ{\,\mathrm{kJ}}
\def\SIs{\,\mathrm{s}}
\def\SIkg{\,\mathrm{kg}}
\def\SIg{\,\mathrm{g}}
\def\SIK{\,\mathrm{K}}
\def\SImmHg{\,\mathrm{mmHg}}
\def\SIPa{\,\mathrm{Pa}}

\def\courseurl{https://github.com/zqhuang/SYSU\_TD}

\def\tpage#1#2{
\begin{frame}
\begin{center}
\begin{Large}
\bch
热学 \\
第#1讲 #2

{\vskip 0.3in}

黄志琦

\ech
\end{Large}
\end{center}

\vskip 0.2in

\bch
教材：《热学》第二版，赵凯华，罗蔚茵，高等教育出版社
\ech

\bch
课件下载
\ech
\courseurl
\end{frame}
}

\title{Lesson 08 Enthalpy}
  \author{}
  \date{}
\begin{document}
\tpage{8}{焓}

\section{Review}

\begin{frame}
\chtitle{上讲需要掌握的内容非常少}
\bch
\bitem
\item{内能是态函数}
\item{热力学第一定律：$\Delta U = A + Q$。}
\item{实际气体等温过程的内能减少量等于内压强做功；普适内能方程$\pfrac UVT = T\pfrac pTV - p$；克拉珀龙方程；范德瓦尔斯气体的内能。}
\item{绝热过程状态方程：$pV^\gamma = \const$ ，其中$\gamma = C_p/C_V = (C_V + \nu R)/C_V$。}
\item{大气的绝热近似$\rightarrow \frac{dT}{dz} \approx -10 \SIK/\SIkm$。}
\item{多方过程气体对外做功$A' = -\frac{\nu R}{n-1}\Delta T$ （等温过程$n=1$则需要另算）。}
\item{多方过程热容$C_n =C_V -\frac{\nu R}{n-1}$ （等温过程$C = \infty$）。}
\eitem
\ech
\end{frame}

\begin{frame}
\chtitle{关于大气垂直温差的补充知识}
\bch
{\small
对水蒸气饱和的空气，做功等于内能改变量加上汽化热：
$$ -pdV = C_V dT + \Lambda d(c\nu) $$
其中$\nu$为空气摩尔数，$c$为空气中水蒸气的摩尔百分比，$\Lambda$为摩尔汽化热。又
$$pdV + Vdp = \nu R dT + \nu \Lambda dc$$
两式相加得到
$$ Vdp = C_p  dT + \nu \Lambda dc $$
再由力学平衡$ dp = -\rho g dz$，即有
$$ \frac{dT}{dz} =-\frac{\rho g V + \nu\Lambda\frac{dc}{dz}}{C_p} = -\frac{\gamma-1}{\gamma}\left(\frac{M^{\rm mol}g}{R}+\frac{\Lambda}{R}\frac{dc}{dz}\right)$$
因为饱和蒸气压随着温度降低而降低（请回忆三相图），在有饱和蒸汽的空气里，气团继续上升会造成水蒸气凝结，$ \frac{dc}{dz} < 0 $，垂直温度梯度就会小于$-10\SIK/\SIs$。
}
\ech
\end{frame}


\begin{frame}
\chtitle{关于空气中声速的补充知识}
\bch

\addfig{1.5}{soundwave.jpg}

{\scriptsize
力学里无敌公式$F = ma$:

$$ -dp dy dz = (\rho dx dy dz) \left(\frac{d\upsilon}{dx/\upsilon}\right) $$
即
$$ dp = - \rho \upsilon d\upsilon $$

然后根据物质流守恒：$ \rho \upsilon = (\rho + d\rho)(\upsilon+d\upsilon)$
忽略高阶小量即$ -\rho d\upsilon = \upsilon d\rho$
代入前面的$dp$表达式得到：
$$ dp = \upsilon ^2 d\rho$$
即
$$\upsilon = \sqrt{\frac{dp}{d\rho}}$$
}
\ech
\end{frame}

\begin{frame}
\chtitle{关于空气中声速的补充知识（续）}
\bch
{\small
空气是热的不良导体，故做绝热近似 
$$p\rho^{-\gamma} = \const$$
即 $$\frac{dp}{d\rho} = \gamma \frac{p}{\rho} =  \frac{\gamma R T}{M^{\rm mol}}$$
其中空气摩尔质量$M^{\rm mol} = 0.0289 \SIkg/\SImol$
算出空气中声速为
$$c_s = \sqrt{\frac{\gamma RT}{M^{\rm mol}}} = 347 \sqrt{\frac{T}{300\SIK}} \SIm/\SIs$$
用无量纲速率来表示就是
$$u_s = \sqrt{\gamma}$$
因此声速和空气分子方均根速率之比为（见课本习题3-21）
$$\sqrt{\frac{\gamma}{3}} = \sqrt{\frac{7}{15}} = 0.683$$
}
\ech
\end{frame}

\begin{frame}
\chtitle{思考题}
\bch

\addfig{1.}{songfen.jpg}

前面关于空气中声速和空气分子的方均根速率的比的计算是有点问题的，你能指出问题在哪里吗？

\ech
\end{frame}


\begin{frame}
\chtitle{本讲内容}
\bch
\bitem
\item{数学准备知识：循环偏微分乘积定理}
\item{焓}
\item{节流过程（等焓的非准静态过程)}
\item{化学反应和生成焓}
\eitem
\ech
\end{frame}

\section{Math}

\begin{frame}
\chtitle{数学准备知识：循环偏微分乘积定理}
\bch
设三个变量$X$, $Y$, $Z$满足某状态方程，则
{\blue 
$$\pfrac XYZ  \pfrac YZX \pfrac ZXY = -1$$
}
当然，我们这是讨论物理问题，可以假设状态方程足够光滑，且不出现偏导数无穷大的情况。

\skipline

\bye不会证的请默默捧起数学书。
\ech
\end{frame}

\begin{frame}
\chtitle{什么，数学书早卖了？}
\bch
{\scriptsize
证明：设状态方程为$S(X, Y, Z) = 0$，取状态方程曲面上的一条等高线($dZ = 0$)，在等高线上有
$$\pfrac SX{YZ} dX + \pfrac SY{ZX} dY = 0$$
又根据偏导数的定义，等高线上的$dX/dY$即$\pfrac XYZ$。故有
$$\pfrac XYZ = -\frac{\pfrac SY{ZX}}{\pfrac SX{ZY}}$$
同理可证
$$\pfrac YZX = -\frac{\pfrac SZ{XY}}{\pfrac SY{ZX}}$$
$$\pfrac ZXY = -\frac{\pfrac SX{YZ}}{\pfrac SZ{XY}}$$
三式相乘即得证
}
\ech
\end{frame}

\begin{frame}
\chtitle{思考题}
\bch
\addfig{1}{songfen2.jpg}

证明气体定体比热容可以写成:

$$
C_V =  - \pfrac UVT \pfrac VTU 
$$

\ech
\end{frame}


\section{Enthalpy}

\begin{frame}
\chtitle{焓}
\bch

内能 $\Rightarrow$ 你房子的价值

焓$\Rightarrow$你房子的价值 + 你房子所占空间的价值

\skipline

自然界很多过程都是等压过程而非等体过程，当体积发生变化时就需要对环境做功以“获得生存空间”。这部分额外做的功做为“无形资产”加到内能上，就是“焓”(enthalpy)。

$$ H \equiv U + pV$$


\ech
\end{frame}

\begin{frame}
\chtitle{思考题}
\bch
\addfig{0.8}{think3.jpg}
焓是态函数吗？
\ech
\end{frame}


\begin{frame}
\chtitle{理想气体的焓只是温度的函数}
\bch
固定摩尔数的理想气体的内能只是温度的函数，故焓
$$H = U(T) + pV = U(T) + \nu R T$$
也只是温度的函数。
\ech
\end{frame}

\begin{frame}
\chtitle{实际气体的焓}
\bch
由于理想气体近似已经给出
$$\pfrac HTp =  C_V+ \nu R > 0$$
实际气体一般只是对上式稍有偏离，并不会影响$\pfrac HTp $的正负符号。

比较不确定的是$H$对$p$的依赖关系，因为理想气体近似给出$\pfrac HpT = 0$，实际气体如果有所偏离的话，是往哪个方向偏离呢？
\ech
\end{frame}

\begin{frame}
\chtitle{$\pfrac HpT$的计算公式}
\bch
{\small
我们上一讲中提到了一个普适内能方程{\blue
$$\pfrac UVT = T\pfrac pTV - p$$}
由此可以得到
$$\pfrac UpT = \pfrac UVT\pfrac VpT= \left(T\pfrac pTV - p\right) \pfrac VpT$$
所以
$$\pfrac HpT = \pfrac UpT  + p\pfrac VpT + V = T\pfrac pTV \pfrac VpT + V$$
最后我们利用循环偏微分乘积定理，得到$\pfrac pTV \pfrac VpT = - \pfrac VTp$，即
{\blue
$$\pfrac HpT = - T\pfrac VTp + V$$
}}
\ech
\end{frame}



\begin{frame}
\chtitle{实际气体的$\pfrac HpT$}
\bch

对范德瓦尔斯气体，我们可以直接算出
$$\pfrac HpT = V\frac{\frac{\nu b}{V} + \frac{2p_U}{p_k}\left(1- \frac{\nu b}{V}\right)}{1 + \frac{2p_U}{p_k}\left(1-\frac{\nu b}{V}\right)}\approx V\left(\frac{\nu b}{V} +\frac{2p_U}{p_k}\right)$$

注意$p_U/p_k<0$，所以上式的符号取决于体积修正和压强修正的相对大小。

\bmini{0.48}
\lfig{1.9}{pfracHpT.jpg}
\emini
\bmini{0.48}
{\small 
实验表明，在常温常压下，氢气和氦气$\pfrac HpT > 0$，而大多数其他气体则有$\pfrac HpT < 0$。

在高温高压下，体积修正是主导项，所有气体均有$\pfrac HpT > 0$。
}
\emini

\ech
\end{frame}

\begin{frame}
\chtitle{思考题}
\bch
\addfig{0.8}{think.jpg}

根据范德瓦尔斯模型，$$\pfrac HpT\approx V\left(\frac{\nu b}{V} + \frac{2p_U}{p_k}\right)$$
可以看成体积修正项($>0$)与压强修正项（$<0$)相互竞争的结果。在常温常压下，对小分子气体，为什么反而结果是$\pfrac HpT > 0$，即体积修正项为主导呢？

\ech
\end{frame}

\begin{frame}
\chtitle{节流过程和焦耳-汤姆孙效应}
\bch

\bmini{0.48}
\addfig{2}{throttling_process.png}
\emini
\bmini{0.48}
用多孔塞把绝热箱隔开，两端分别施以恒定压强$p_1$, $p_2$ （$p_1>p_2$）。
\emini

{\small

假设一开始气体都在左边，体积为$V_1$，内能为$U_1$。经过一段时间，左边的气体都通过多孔塞被压到了右边，体积和内能变为$V_2$, $U_2$。

那么由于左边活塞对气体做功$p_1V_1$，右边活塞对气体做功$-p_2V_2$，所以气体内能变化为

$$U_2 - U_1  = p_1V_1 - p_2V_2$$

即
$$ U_1 + p_1V_1 = U_2 + p_2V_2 $$
{\blue \bf 节流过程是等焓过程} （注意：节流过程一般不是准静态过程）

}
\ech
\end{frame}


\begin{frame}
\chtitle{焦耳-汤姆孙效应}
\bch
{\small
我们可以根据焦耳-汤姆孙系数$\alpha \equiv \pfrac TpH$来计算的温度的升降。

根据循环偏微分乘积定理：
$$\pfrac TpH \pfrac pHT \pfrac HTp = -1$$
和$\pfrac HTp > 0$，就知道$\pfrac TpH$和$\pfrac HpT$符号相反。

也就是说，在常温常压下，对氢气和氦气$\pfrac TpH < 0$， 对大多数其他气体$\pfrac TpH > 0$。

因为节流过程是降压过程，所以对{\bf 大部分气体（空气，氮气，氧气等）节流过程是降温过程（正节流效应），而对氢气和氦气是升温过程（负节流效应）}。这些效应统称{\bf 焦耳-汤姆孙效应}。
}
\ech
\end{frame}


\begin{frame}
\chtitle{节流液化气体}
\bch
{\small
根据范德瓦尔斯状态方程的定性分析，在高温高压下，气体呈现负节流效应，在低温低压下，气体呈现正节流效应。那么只要把气体降温到某个温度之下使得节流效应为正，则节流效应可以对气体无限降温使它最终液化。

\skiplines

请参考教材图3-15的节流液化装置。

\skiplines

\bitem
\item{实际气体在极低温时未必符合范德瓦尔斯模型，具体例子可见教材140页图3-13，氢气在极低温时也有一部分负节流效应区域。这样氢气的节流液化过程就会有一定的参数限制。}
\item{自由膨胀也能给气体降温。柯林斯液化机就是结合自由膨胀和节流两种方法。}
\eitem

}
\ech
\end{frame}

\section{Enthalpy of Reaction}


\begin{frame}
\chtitle{化学反应热和生成焓}
\bch
定义化学反应热 = 外界给系统的热量

\skipline

\bitem
\item{定体化学反应 $\Rightarrow$ 内能变化 = 反应热}
\item{定压化学反应 $\Rightarrow$ 焓变化 = 反应热}
\eitem

\skipline

定压反应的反应热又称为{\blue 反应焓}。
\ech
\end{frame}

\begin{frame}
\chtitle{焓的参考点：标准生成焓}
\bch
焓和势能一样，可以取任意参考点，{\blue 习惯上定义 $1\atm$, $25\cdeg$的纯元素的焓为零}。

\skipline

{\blue $1\atm$，$25\cdeg$下化合物的生成焓定义为化合物的标准生成焓(用$\stdHf$表示)}。

\skipline

{\blue $1\atm$，$25\cdeg$下，

 反应焓 = 合成物的标准生成焓 - 反应物的标准生成焓}

\skiplines

标准生成焓可以通过查表获得。见教材145页表3-3。

\ech
\end{frame}

\begin{frame}
\chtitle{反应焓计算举例}
\bch
\bex
教材习题3-23
\eex

\skipline

记给的三个燃烧反应为$A, B, C$。则 $2A + \frac{3}{2}B - \frac{1}{2}C$就等价于所需要的纯元素合成反应。故
{\small
\bea
\stdHf({\rm C_2H_6}) &=&\left[ 2\times (-393.5) + \frac{3}{2}\times (-571.7)-\frac{1}{2}\times(-1560)\right]\SIkJ/\SImol\newl
&=&-864.6\SIkJ/\SImol
\eea
}
\ech
\end{frame}

\begin{frame}
\chtitle{第九周作业（序号接第八周）}
\bch
\bitem
\item[22]{证明气体的定压比热容$C_p$可以写成
$$C_p = - \pfrac pTH \pfrac HpT $$}
\item[23]{教材习题3-14}
\item[24]{教材习题3-24}
\eitem
\ech
\end{frame}
\end{document}
